\section{مدل تفصیلی ظرفیت \lr{Gate}}
\subsection{سهم \lr{Intrinsic} و وابستگی صفحه پایین به ناحیه کاری}
در مدل تفصیلی بخش \secnum{۲.۳.۲}، ظرفیت \lr{Gate} شامل سهم روی کانال و سهم
هم‌پوشانی با \lr{Source} و \lr{Drain} است. سهم روی کانال، \lr{Intrinsic}
نام دارد. ظرفیت هندسی مرجع را به صورت
\begin{equation}
 C_0=WLC_{\mathrm{ox}}
 \label{eq:cv-czero}
\end{equation}
تعریف می‌کنیم. این کمیت مرجع برای مقایسه ناحیه‌هاست. تقسیم ظرفیت میان پایانه‌ها
به محل بار پاسخ‌دهنده به \lr{Gate} بستگی دارد
\cite[صفحه ۷۰، شکل \secnum{۲.۹}]{weste}.

در \lr{Cutoff}، کانال \lr{Strong Inversion} وجود ندارد. تغییر بار \lr{Gate}
عمدتاً با تغییر بار \lr{Body} جبران می‌شود. سهم مربوط، \(C_{GB,i}\) است.
در \lr{Accumulation}، حامل‌های غالب نزدیک سطح‌اند و این ظرفیت تقریباً
\(C_0\) است. در \lr{Depletion}، ناحیه کم‌حامل مانند فاصله اضافی میان بارهای
پاسخ‌دهنده عمل می‌کند و ظرفیت مؤثر کاهش می‌یابد.

برای درک مدل سری، \(C_{\mathrm{dep}}\) را ظرفیت کل ناحیه \lr{Depletion}
زیر \lr{Gate} می‌نامیم. چون در دو بخش سری، بزرگی تغییر بار یکسان و تغییرهای
ولتاژ جمع‌شونده‌اند، داریم
\begin{equation}
 \frac1{C_{GB,i}}\approx\frac1{C_0}+\frac1{C_{\mathrm{dep}}}
 \quad\Longrightarrow\quad
 C_{GB,i}\approx\frac{C_0C_{\mathrm{dep}}}{C_0+C_{\mathrm{dep}}}.
 \label{eq:cv-depletion-series}
\end{equation}
این توضیح تکمیلی، کاهش ظرفیت نشان‌داده‌شده در شکل \secnum{۲.۹} کتاب را بیان
می‌کند. اگر هر دو ظرفیت سری برابر \qty{6}{\femto\farad} باشند، حاصل
\(36/12=\qty{3}{\femto\farad}\) است.
افزایش پهنای ناحیه کم‌حامل، ظرفیت آن را کم می‌کند و ظرفیت سری نیز کاهش می‌یابد.

در \lr{Linear} با \(V_{GS}>V_{tn}\) و \(V_{DS}\) نزدیک صفر، کانال به هر دو
\lr{Source} و \lr{Drain} متصل است. مدل تقریبی کتاب سهم را برابر تقسیم می‌کند
\begin{equation}
 C_{GS,i}\approx\frac{C_0}{2},\qquad
 C_{GD,i}\approx\frac{C_0}{2},\qquad C_{GB,i}\approx0.
 \label{eq:cv-linear-partition}
\end{equation}
با افزایش \(V_{DS}\) در همان ناحیه، بار کانال نزدیک \lr{Drain} کمتر می‌شود.
سهم منتسب به \lr{Source} افزایش و سهم منتسب به \lr{Drain} کاهش می‌یابد.
جدول ساده کتاب، مقادیر نزدیک \(V_{DS}=0\) را برای ناحیه \lr{Linear} نمایندگی می‌کند.

در \lr{Saturation}، انتهای کانال از \lr{Drain} جدا شده و در این تقریب،
\(C_{GD,i}\approx0\) و \(C_{GB,i}\approx0\) هستند. ظرفیت مؤثر ذاتی به
\lr{Source} نسبت داده می‌شود و مقدار آن
\begin{equation}
 C_{GS,i}\approx\frac23C_0
 \label{eq:cv-sat-cap}
\end{equation}
است \cite[صفحه ۷۰ و جدول \secnum{۲.۱}]{weste}.
در مدل پیوسته، تغییر مؤلفه‌ها با نزدیک‌شدن به مرز نرم است.
جدول قطعه‌ای، این تغییر را برای تخمین دستی خلاصه می‌کند.

\clearpage
\subsection{استخراج ضریب دوسوم از بار کانال}
برای فهم ضریب \(2/3\)، از مدل بار موضعی بخش \secnum{۲.۲} استفاده می‌کنیم.
در این مدل، \(x\) فاصله از \lr{Source} و \(V(x)\) ولتاژ محلی کانال نسبت به
\lr{Source} است. اضافه‌ولتاژ گیت را با شرط مثبت‌بودن آن تعریف می‌کنیم.
\[
 V_{GT}=V_{GS}-V_{tn},\qquad V_{GT}>0.
\]
بزرگی بار متحرک قطعه‌ای به طول \(dx\)
برابر \(WC_{\mathrm{ox}}[V_{GT}-V(x)]dx\) است.
رابطه جریان محلی با \lr{Mobility} و آستانه ثابت، پس از انتگرال از صفر تا \(x\)، می‌دهد
\begin{equation}
 I_{DS}x=\mu_nWC_{\mathrm{ox}}
 \left[V_{GT}V(x)-\frac{V(x)^2}{2}\right].
 \label{eq:cv-potential-profile}
\end{equation}
در اشباع بلندکانال، \(I_{DS}=\mu_nC_{\mathrm{ox}}(W/2L)V_{GT}^2\) است.
با جای‌گذاری و تقسیم دو طرف بر \(\mu_nWC_{\mathrm{ox}}V_{GT}^2/2\)، نتیجه می‌شود
\begin{equation}
 \frac{x}{L}=2\frac{V(x)}{V_{GT}}-
 \left(\frac{V(x)}{V_{GT}}\right)^2,
 \qquad
 1-\frac{x}{L}=\left(1-\frac{V(x)}{V_{GT}}\right)^2.
 \label{eq:cv-square-profile}
\end{equation}
چون \(0\le V(x)\le V_{GT}\)، ریشه نامنفی را می‌گیریم، پس
\begin{equation}
 V_{GT}-V(x)=V_{GT}\sqrt{1-\frac{x}{L}}.
 \label{eq:cv-root-profile}
\end{equation}
بار متحرک در نزدیکی \lr{Source} بیشینه و در انتهای مدل کانال صفر است.
جمع بار همه قطعه‌ها، با تغییر متغیر بی‌بعد \(u=x/L\)، می‌دهد
\begin{align}
 Q_{\mathrm{ch}}
 &=WC_{\mathrm{ox}}V_{GT}\int_0^L\sqrt{1-\frac{x}{L}}\,dx\notag\\
 &=WLC_{\mathrm{ox}}V_{GT}\int_0^1\sqrt{1-u}\,du
 =\frac23 C_0V_{GT}.
 \label{eq:cv-sat-charge}
\end{align}
انتگرال آخر مساحت زیر منحنی \(\sqrt{1-u}\) است. برای محاسبه مرحله‌ای،
پادمشتق آن \(-\tfrac23(1-u)^{3/2}\) است. مقدار در انتهای بازه صفر و در
ابتدای بازه \(-2/3\) می‌شود. تفاضل دو مقدار، \(2/3\) است.
عامل \(2/3\) از شکل توزیع بار در کانال حاصل می‌شود.

در این مدل، بخش متغیر بار \lr{Gate} با بزرگی بار کانال متناظر است.
با ثابت‌گرفتن آستانه و سایر بایاس‌های لازم در تقریب ذاتی، تغییر \(V_{GS}\)
همان تغییر \(V_{GT}\) است. بنابراین
\begin{equation}
 C_{GS,i}\approx\frac{dQ_G}{dV_{GS}}
 =\frac{d[(2/3)C_0(V_{GS}-V_{tn})]}{dV_{GS}}
 =\frac23 C_0.
 \label{eq:cv-sat-derivative}
\end{equation}
این استخراج تکمیلی، مقدار جدول کتاب را در فرض‌های ایدئال توضیح می‌دهد.
مدل دقیق‌تر مبتنی بر بار پایانه‌ها، تغییرات بار و جریان را با روش سازگار مدل
فناوری محاسبه می‌کند. کتاب در پایان این زیربخش به \lr{Charge-based Model} اشاره دارد.

\subsection{جدول ناحیه‌ها و تفسیر \lr{Cutoff}}
جدول \ref{tab:cv-intrinsic} مؤلفه‌های ذاتی را یکجا نشان می‌دهد.
مؤلفه‌های \lr{Overlap} در این جدول به صورت جداگانه در نظر گرفته خواهند شد.
\begin{table}[ht]
\centering
\caption{مدل تقریبی ظرفیت ذاتی، با روشن‌کردن شرط‌های جدول \secnum{۲.۱} کتاب.}
\label{tab:cv-intrinsic}
\begin{tabular}{cccc}
\toprule
کمیت & \lr{Cutoff} & \lr{Linear} با \(V_{DS}\approx0\) & \lr{Saturation}\\
\midrule
\(C_{GB,i}\) & \(\le C_0\) & \(0\) & \(0\)\\
\(C_{GS,i}\) & \(0\) & \(C_0/2\) & \(2C_0/3\)\\
\(C_{GD,i}\) & \(0\) & \(C_0/2\) & \(0\)\\
\(C_{gc}=C_{GB,i}+C_{GS,i}+C_{GD,i}\) & \(\le C_0\) & \(C_0\) & \(2C_0/3\)\\
\bottomrule
\end{tabular}
\end{table}
در سطر آخر جدول کتاب، برای \lr{Cutoff} مقدار ساده‌شده \(C_0\) آمده است،
درحالی‌که سطر \(C_{GB}\) و متن صفحه ۷۰، کاهش ظرفیت در \lr{Depletion} را بیان
می‌کنند. مقدار \(C_0\) با تقریب \lr{Accumulation} سازگار است.
در این جزوه، علامت \(\le\) در سطر مجموع حفظ شده تا وابستگی بایاس آشکار بماند.
کتاب هشدار می‌دهد که استفاده همیشگی از \(C_0\)، ظرفیت افزاره نزدیک زیر آستانه
را به‌طور قابل‌توجهی بیش‌برآورد می‌کند \cite[صفحه ۷۲]{weste}.

در \lr{Cutoff}، جریان هدایت کانال در مدل مرتبه اول صفر است، اما با تغییر
ولتاژ، بار \lr{Gate} و ناحیه‌های پیوندی می‌تواند تغییر کند.
به همین دلیل، مدار خاموش نیز در تحلیل گذرا ظرفیت دارد.

\subsection{\lr{Overlap}، \lr{Fringing} و ظرفیت کل}
در افزاره واقعی، \lr{Gate} بخشی از ناحیه‌های \lr{Source} و \lr{Drain} را
می‌پوشاند و میدان‌های لبه‌ای، یا \lr{Fringing Fields}، نیز روی آن‌ها ختم می‌شوند.
این سهم‌ها در مدل کتاب به ظرفیت‌های \lr{Overlap} افزوده می‌شوند.
ضریب بر واحد عرض آن‌ها را \(c_{GS,\mathrm{ov}}\) و
\(c_{GD,\mathrm{ov}}\) می‌نامیم. این دو نماد متناظر با ضرایب
\(C_{gsol}\) و \(C_{gdol}\) در رابطه \secnum{۲.۱۵} کتاب هستند.
ظرفیت کل هر سهم چنین است
\begin{equation}
 C_{GS,\mathrm{ov}}=c_{GS,\mathrm{ov}}W,\qquad
 C_{GD,\mathrm{ov}}=c_{GD,\mathrm{ov}}W.
 \label{eq:cv-overlap}
\end{equation}
کتاب مقادیر معمول آن مثال‌ها را حدود \num{0.2} تا
\qty{0.4}{\femto\farad\per\micro\metre} بیان می‌کند
\cite[صفحه ۷۱، شکل \secnum{۲.۱۰}]{weste}.
در مدل ضرایب ثابت این بخش، سهم \lr{Overlap} با عرض رشد می‌کند.

با قرارداد این جزوه، جمع سهم ذاتی و هم‌پوشانی به صورت زیر است
\begin{equation}
 C_{GS}=C_{GS,i}+C_{GS,\mathrm{ov}},\qquad
 C_{GD}=C_{GD,i}+C_{GD,\mathrm{ov}},\qquad
 C_{GB}\approx C_{GB,i}.
 \label{eq:cv-total-components}
\end{equation}
در اشباع، \(C_{GD,i}\) در جدول ایدئال صفر است و \(C_{GD}\) می‌تواند سهم
\lr{Overlap} داشته باشد. اگر برای هر دو سمت ضریب یکسان \(c_{\mathrm{ov}}\)
به کار رود، تقریب کلی کتاب برای تخمین تأخیر
\(C_g\approx C_0+2c_{\mathrm{ov}}W\) است.
برای تحلیل ناحیه‌ای با پایانه‌های ثابت، سهم ذاتی همان ناحیه جایگزین \(C_0\) می‌شود.

مثال فرضی \(C_0=\qty{6}{\femto\farad}\)، \(W=\qty{2}{\micro\metre}\)
و \(c_{\mathrm{ov}}=\qty{0.3}{\femto\farad\per\micro\metre}\) را در نظر بگیرید.
هر سهم \lr{Overlap} برابر \qty{0.6}{\femto\farad} می‌شود.
در \lr{Linear} کم‌ولتاژ، دو ظرفیت برابرند.
\[
 C_{GS}=C_{GD}=3+0.6=\qty{3.6}{\femto\farad}.
\]
مجموع ظرفیت‌ها \qty{7.2}{\femto\farad} است. در اشباع، داریم
\[
 C_{GS}=4+0.6=\qty{4.6}{\femto\farad},\qquad
 C_{GD}=\qty{0.6}{\femto\farad}.
\]
مجموع ظرفیت‌های این حالت \qty{5.2}{\femto\farad} است.
در \lr{Accumulation}، مجموع تقریباً \(6+1.2=\qty{7.2}{\femto\farad}\) می‌شود.
در \lr{Depletion} فرضی با \(C_{\mathrm{dep}}=\qty{3}{\femto\farad}\)،
ظرفیت ذاتی از مدل سری به دست می‌آید.
\[
 C_{GB,i}=\frac{6(3)}{6+3}=\qty{2}{\femto\farad}.
\]
پس مجموع ظرفیت‌ها \qty{3.2}{\femto\farad} است.

\subsection{حرکت \lr{Source} و \lr{Drain} و ظرفیت مؤثر ورودی}
برای یک سهم ثابت \(C_{GX}\) میان \lr{Gate} و پایانه \(X\)، تغییر بار
\lr{Gate} برابر است با
\begin{equation}
 \Delta Q_G=C_{GX}(\Delta V_G-\Delta V_X).
 \label{eq:cv-moving-terminal}
\end{equation}
اگر \(\Delta V_G\ne0\)، با تقسیم بر تغییر ولتاژ ورودی، سهم مؤثر آن به دست می‌آید
\begin{equation}
 C_{\mathrm{eff},G}=C_{GX}\left(1-\frac{\Delta V_X}{\Delta V_G}\right).
 \label{eq:cv-miller-simple}
\end{equation}
در مدل خازن ثابت، برای پایانه ثابت، ضریب یک است.
اگر دو پایانه به یک اندازه و هم‌جهت حرکت کنند، اختلاف ولتاژ خازن تغییر نمی‌کند
و سهم جابه‌جایی بار آن در این گذار صفر است.
اگر هم‌اندازه و خلاف‌جهت حرکت کنند، ضریب دو می‌شود.
این تحلیل ساده، منشأ اثر حرکت پایانه دوم را نشان می‌دهد.

برای مثال، خازن \qty{1}{\femto\farad} میان دو گره را در نظر بگیرید.
\lr{Gate} از صفر به \qty{1}{\volt} می‌رود. اگر گره دیگر صفر بماند، تغییر بار
\qty{1}{\femto\coulomb} است. اگر گره دیگر از \qty{1}{\volt} به صفر برود،
تغییر اختلاف ولتاژ \qty{2}{\volt} و تغییر بار \qty{2}{\femto\coulomb} است.
اگر گره دیگر نیز از صفر به \qty{1}{\volt} برود، تغییر بار همین خازن صفر می‌شود.

کتاب در شکل \secnum{۲.۱۱} هفت ترکیب حرکت \lr{Source} و \lr{Drain} را برای
یک فرآیند \qty{0.35}{\micro\metre} گزارش می‌کند.
جدول \ref{tab:cv-data-dependent} داده‌های آن شکل را بازنویسی می‌کند.
در همه حالت‌ها، \lr{Gate} از صفر به یک می‌رود و \(C_g/C_0\) نسبت مؤثر گزارش‌شده
همان شکل است. این اعداد به افزاره و مسیر گذار مورد بررسی وابسته‌اند.
\begin{table}[ht]
\centering
\caption{وابستگی ظرفیت مؤثر \lr{Gate} به حرکت پایانه‌ها در شکل
\secnum{۲.۱۱} کتاب. پیکان، جهت گذار و عدد، مقدار منطقی ثابت را نشان می‌دهد \cite{weste}.}
\label{tab:cv-data-dependent}
\begin{tabular}{cccc}
\toprule
حالت شکل کتاب & \lr{Drain} & \lr{Source} & \(C_g/C_0\)\\
\midrule
۱ & \(\downarrow\) & \(\downarrow\) & \(1.3\)\\
۲ & \(\downarrow\) & \(0\) & \(1.1\)\\
۳ & \(0\) & \(0\) & \(1.0\)\\
۴ & \(1\) & \(0\) & \(0.80\)\\
۵ & \(1\) & \(1\) & \(0.42\)\\
۶ & \(1\) & \(\uparrow\) & \(0.31\)\\
۷ & \(\uparrow\) & \(\uparrow\) & \(0.13\)\\
\bottomrule
\end{tabular}
\end{table}
در ترانزیستور، مؤلفه‌های ظرفیت طی گذار نیز تغییر می‌کنند.
برای محاسبه دقیق‌تر باید بار پایانه‌ها در مسیر واقعی بایاس محاسبه شود.
مثال خازن ثابت، وابستگی به جهت حرکت را توضیح می‌دهد و اعداد جدول از مدل افزاره
و گذارِ منبع آمده‌اند. کتاب برای تخمین تأخیر، استفاده از ظرفیت مؤثر استخراج‌شده
از شبیه‌سازی را نیز مطرح می‌کند \cite[صفحات ۷۱ و ۷۲]{weste}.

\subsection{تمرین‌های پاسخ‌دار \lr{Gate Capacitance}}
\textbf{تمرین ۴.} برای \(C_0=\qty{9}{\femto\farad}\)، مؤلفه‌های ذاتی در
\lr{Linear} با \(V_{DS}\approx0\) و در اشباع چیست؟ در حالت اول،
\(C_{GS,i}=C_{GD,i}=\qty{4.5}{\femto\farad}\) و \(C_{GB,i}=0\) است.
در اشباع، \(C_{GS,i}=\qty{6}{\femto\farad}\) و دو سهم دیگر صفرند.

\textbf{تمرین ۵.} در تمرین قبل، \(W=\qty{3}{\micro\metre}\) و ضریب هر
\lr{Overlap} برابر \qty{0.2}{\femto\farad\per\micro\metre} باشد.
در اشباع، ظرفیت‌های کل و مجموعشان چقدر است؟ هر هم‌پوشانی
\qty{0.6}{\femto\farad} است. پس \(C_{GS}=\qty{6.6}{\femto\farad}\)،
\(C_{GD}=\qty{0.6}{\femto\farad}\) و مجموع \qty{7.2}{\femto\farad} می‌شود.

\textbf{تمرین ۶.} خازن ثابت \qty{2}{\femto\farad} به \lr{Gate} با تغییر
\qty{1}{\volt} متصل است. پایانه دیگر \qty{0.25}{\volt} هم‌جهت تغییر می‌کند.
ظرفیت مؤثر از دید \lr{Gate} چیست؟ از \eqref{eq:cv-miller-simple}، مقدار
\(2(1-0.25)=\qty{1.5}{\femto\farad}\) و تغییر بار
\qty{1.5}{\femto\coulomb} می‌شود.

\textbf{تمرین ۷.} در مدل سری \lr{Depletion}، \(C_0=\qty{8}{\femto\farad}\)
و \(C_{\mathrm{dep}}=\qty{8}{\femto\farad}\) باشد. اگر ظرفیت دوم نصف شود،
\(C_{GB,i}\) چه تغییری می‌کند؟ ابتدا \(64/16=\qty{4}{\femto\farad}\) است.
پس از تغییر، \(8(4)/(8+4)\approx\qty{2.67}{\femto\farad}\) می‌شود.
\clearpage
